%
% 12-algebraisch.tex -- Algebraische Erweiterungen von Q
%
% (c) 2026 Prof Dr Andreas Müller
%
\subsection{Algebraische Erweiterungen von $\mathbb{Q}$}
Die Vervollständigung der der rationalen Zahlen zu den reellen
Zahlen stellt sicher, dass jede Cauchy-Folge einen Grenzwert hat.
Sie führt aber zu einem neuen Problem, welches ein Hindernis für
Computeralgebrasysteme wird.
\index{Computeralgebra}%
Da reelle Zahlen nur durch ihre Approximationen definiert sind,
lässt sich damit grundsätzlich nur dann exakt rechnen, wenn auch
noch eine alternative Charakterisierung zur Verfügung steht.
Mit der Kreiszahl $\pi$ kann man nur desshalb exakt rechnen, weil
sie die kleinste positive Nullstelle der Sinusfunktion ist.

Für die Quadratwurzel $\!\sqrt{2}$, die früher als erstes Beispiel
\index{Quadratwurzel}%
einer irrationalen Zahl in den Vordergrund gestellt wurde, gibt
es eine solche alternative Charakterisierung.
Beim algebraischen Rechnen mit $\!\sqrt{2}$ verwendet man fast
instinktiv die Beziehung $(\!\sqrt{2})^2=2$, zum Beispiel
beim Ausmultiplizieren von
\[
(\!\sqrt{2}+1)^2
=
(\!\sqrt{2})^2
+2\!\sqrt{2}\cdot 1 
+ 1^2
=
2 + 2\!\sqrt{2} + 1
=
3+2\!\sqrt{2}.
\]

Wir wollen diese Beobachtung noch etwas expliziter machen und
bezeichnen mit $\alpha$ ein Objekt, für welches die Regeln
der Algebra ebenso gelten wie die Identität $\alpha^2-2=0$.
Durch die Ersetzung $\alpha^2 = 2$ kann jede Potenz von $\alpha$
auf entweder ein Vielfaches von $\alpha$ oder auf eine ganze Zahl
reduziert werden, es gilt nämlich
\[
\alpha^n
=
\begin{cases}
2^{\frac{n}2}         &\qquad\text{$n$ gerade}\\
2^{\frac{n-1}2}\alpha &\qquad\text{$n$ ungerade.}
\end{cases}
\]
Für Linearkombinationen $a+b\alpha$ und $c+d\alpha$ lassen sich
jetzt einfache Rechenregeln aufstellen:
\begin{align*}
(a+b\alpha)\pm(c+d\alpha) &= (a\pm b) + (c\pm d)\alpha
\\
(a+b\alpha)(c+d\alpha) &= (ac + 2bd) + (ad+bc)\alpha.
\end{align*}
Als Spezialfall der Produktregel findet man
\[
(a+b\alpha)(a-b\alpha)
=
a^2 - 2 b^2,
\]
daher heisst $a-b\alpha$ auch das zu $a+b\alpha$ \emph{konjugierte}
Element.
\index{konjugiert}%
Durch Erweitern mit $(c-d\alpha)$ kann man sogar den Kehrwert
\begin{align*}
\frac{1}{c+d\alpha}
&=
\frac{1}{c+d\alpha}
\cdot
\frac{c-d\alpha}{c-d\alpha}
=
\frac{c-d\alpha}{c^2-2d^2}
=
\frac{c}{c^2-2d^2} - \frac{d}{c^2-2d^2}\alpha
\intertext{und damit auch den Quotienten}
\frac{a+b\alpha}{c+d\alpha}
&=
\frac{ac-2bd}{c^2-2d^2} - \frac{ad-bc}{c^2-2d^2}\alpha
\end{align*}
definieren.

Alle vier Grundoperationen lassen sich also für alle Zahlen der Form 
$a+b\alpha$ ausschliesslich durch Verwendung der arithmetischen Gesetze
und der Identität $\alpha^2-2$ definieren und sie ergeben immer wieder
Zahlen dieser Form.
Es entsteht somit eine neue Menge
\[
\mathbb{Q}(\alpha)
=
\{a+b\alpha\mid a,b\in\mathbb{Q}\},
\]
in der die gleichen Rechengesetze gelten wie in $\mathbb{Q}$.
Die Menge $\mathbb{Q}$ ist eine echte Teilmenge, $\mathbb{Q}(\alpha)$
ist also eine echte Erweiterung von $\mathbb{Q}$.
In $\mathbb{Q}(\alpha)$ lässt sich jetzt die Gleichung $X^2-2=0$ lösen.
Es ist nämlich
\[
(X+\alpha)(X-\alpha)
=
X^2-\alpha^2
=
X^2 - 2
=
0,
\]
die beiden Lösungen der Gleichung $X^2-2=0$ sind also $\pm\alpha$.

Die Konstruktion lässt sich noch etwas verallgemeinern.
Ein \emph{normiertes Polynom} ist ein Polynom dessen führender Koeffizient
\index{normiertes Polynom}%
\index{Polynom!normiert}%
$1$ ist.
Verwendet man statt des Polynoms $X^2-2$ ein beliebiges normiertes
Polynom $m(X)$ aus der Menge
\[
\mathbb{Q}[X]
=
\{
X^n+m_{n-1}X^{n-1}+\dots+m_1X+m_0 
\mid
m_0,\dots,m_{n-1}\in\mathbb{Q}
\}
\]
der Polynome mit rationalen Koeffizienten ersetzt wird\footnote{Im Folgenden
werden wir die Variable eines Polynoms mit grossen Buchstaben bezeichnen.
Sie wird damit klar von den Werten unterschieden, die man anstelle der
Variable ein setzen kann.
Die Variable $X$ kann man auch als ein transzendentes Element über dem
Koeffizientenkörper betrachten.
}.
Wieder versucht man ein neues Objekt $\alpha$ durch die Eigenschaft
$m(\alpha)=0$ zu charakterisieren.
Dies führt jedoch auf Schwierigkeiten, wenn sich $m$ als Produkt
$m=ab$ von Polynomen $a(X),b(X)\in\mathbb{Q}[X]$ schreiben lässt,
denn die beiden Faktoren $a(X)$ und $b(X)$ könnten Nullstellen mit
widersprüchlichen Eigenschaften definieren.

\begin{definition}[irreduzibles Polynom]
Das Polynom $m(X)\in\mathbb{Q}[X]$ heisst \emph{irreduzibel}, wenn es keine
\index{irreduzibel}%
Polynome vom Grad $>0$ gibt, die $m$ als Produkt haben.
\end{definition}

Man beachte, dass die Eigenschaft, irreduzibel zu sein, vom
Koeffizientenkörper abhängt.
Das Polynom $X^2-2$ ist irreduzibel in $\mathbb{Q}[X]$, es lässt sich
nicht in zwei Faktoren zerlegen.
Die einzig mögliche Faktorisierung wäre $X^2-2=(X+\!\sqrt{2})(X-\!\sqrt{2})$,
aber die Faktoren $X+\!\sqrt{2}$ und $X-\!\sqrt{2}$ haben nicht
alle Koeffizienten in $\mathbb{Q}$.
Das Polynom $X^2-2\in\mathbb{R}[X]$ ist also nicht irreduzibel als Polynom
mit reellen Koeffizienten.

In einem irreduziblen Polynom $m(X)$ kann der Koeffizient $m_0$ nicht
verschwinden, denn dann liesse sich
\[
m
=
(X^{n-1} + m_{n-1}X^{n-2}+\dots + m_1)X
\]
als Produkt zweier Polynome vom Grad $n-1$ bzw.~1 schreiben.

Mit einem irreduziblen Polynom $m(X)$ definieren wir jetzt wieder das
Objekt $\alpha$ mit der Eigenschaft $m(\alpha)=0$.
$m$ heisst auch das \emph{Minimalpolynom} von $\alpha$.
\index{Minimalpolynom}%
Ausgeschrieben bedeutet dies
\[
m(\alpha)
=
\alpha^n
+
m_{n-1}\alpha^{n-1}
+
\dots
+
m_1\alpha
+
m_0
=
0
\quad\Rightarrow\quad
\alpha^n
=
-m_{n-1}\alpha^{n-1}
-
\dots
-
m_1\alpha
-
m_0.
\]
Man kann also jede Potenz grösser oder gleich $n$ von $\alpha$ durch
eine Linearkombination von Potenzen $\alpha^k$ mit $0\le k< n$ 
ausdrücken.

\begin{definition}[algebraische Erweiterung]
Sei $m(X)\in\mathbb{Q}[X]$ ein irreduzibles Polynom vom Grad $n$
der Form
\[
m(X)
=
X^n + m_{n-1}X^{n-1} + \dots + m_1X + m_0.
\]
Die Menge
\begin{equation}
\mathbb{Q}(\alpha)
=
\{
a_0 + a_1\alpha + \dots + a_{n-1}\alpha^{n-1}
\mid
a_k\in\mathbb{Q}\text{ mit } 0\le k<n
\}
\label{buch:komplex:zahlen:def:eqn:Qa}
\end{equation}
heisst die \emph{algebraische Erweiterung} von $\mathbb{Q}$ um $\alpha$.
\index{Erweiterung, algebraisch}%
\index{algebraische Erweiterung}%
\end{definition}

Wieder sollen die Rechengesetze durch Übertragung der gewohnten
Regeln entstehen, jedoch müssen alle höheren Potenzen von $\alpha$
jeweils wieder auf Potenzen $\alpha^k$, $0\le k<n$, reduziert werden.
Da der Koeffizient $m_0$ des irreduziblen Polynoms $m$ nicht verschwindet,
kann man die Gleichung $m(\alpha) = 0$ mit $\alpha^{-1}$ multiplizieren
und nach $\alpha^{-1}$ auflösen:
\begin{align*}
m(\alpha)\cdot \alpha^{-1}
=
0
&=
\alpha^{n-1}
+
m_{n-1}\alpha^{n-2}
+
\dots
+
m_1
+
m_0\alpha^{-1}
\\
\alpha^{-1}
&=
-\frac{1}{m_0}
(
\alpha^{n-1}
+
m_{n-1}\alpha^{n-2}
+
\dots
+
m_1
)
\\
&=
-\frac{m_1}{m_0}
-\dots
-\frac{m_{n-1}}{m_0}\alpha^{n-2}
-\frac{1}{m_0}\alpha^{n-1}
\in
\mathbb{Q}(\alpha).
\end{align*}
Das reziproke Element ist also immer auch in $\mathbb{Q}(\alpha)$.
\index{reziprok}%

\begin{beispiel}
\label{buch:komplex:zahlen:bsp:kubikwurzel}
Die Gleichung $X^3-2=0$ lässt sich in $\mathbb{Q}$ nicht lösen.
Das Polynom $m(X)=X^3-2$ ist irreduzibel als Polynom in $\mathbb{Q}[X]$.
Die Erweiterung $\mathbb{Q}(\alpha)$ um das Objekt $\alpha$ mit
\index{Kubikwurzel}%
der Eigenschaft $\alpha^3-2=0$ ist also die Menge
\[
\mathbb{Q}(\alpha)
=
\{
a_0+a_1\alpha + a_2\alpha^2
\mid
a_0,a_1,a_2\in\mathbb{Q}
\}.
\]
Die Additionsregeln sind nicht weiter anspruchsvoll, wir berechnen
daher nur das Produkt von zwei Elementen $a,b\in\mathbb{Q}(\alpha)$:
\begin{align*}
ab
&=
(a_0+a_1\alpha+a_2\alpha^2)
(b_0+b_1\alpha+b_2\alpha^2)
\\
&=
a_0b_0
+
(a_0b_1+a_1b_0)\alpha
+
(a_0b_2+a_1b_1+a_2b_0)\alpha^2
+
(a_1b_2+a_2b_1){\color{darkred}\alpha^3}
+
a_2b_2{\color{darkred}\alpha^4}.
\intertext{Die Potenzen $\alpha^3$ und $\alpha^4$ müssen durch Anwendung der
Reduktionsregel $\alpha^3=2$ auf}
&=
a_0b_0
+
(a_0b_1+a_1b_0)\alpha
+
(a_0b_2+a_1b_1+a_2b_0)\alpha^2
+
(a_1b_2+a_2b_1){\color{darkred}\cdot 2}
+
a_2b_2{\color{darkred} \cdot 2\alpha}
\intertext{reduziert werden.
Das Produkt ist daher}
ab
&=
(a_0b_0 + {\color{darkred}2}a_1b_2+{\color{darkred}2}a_2b_1)
+
(a_0b_1+a_1b_0 + {\color{darkred}2}a_2b_2)\alpha
+
(a_0b_2+a_1b_1+a_2b_0)\alpha^2
\in\mathbb{Q}(\alpha).
\end{align*}
Das zu $\alpha$ reziproke Element ist
\[
\alpha^{-1}
=
\frac{1}{2}\alpha^2,
\]
tatsächlich ist
\[
\alpha\cdot \alpha^{-1}
=
\alpha
\cdot
\frac{1}{2}\alpha^2
=
\frac{\alpha^3}{2}
=
\frac{2}{2}
=
1.
\]
In $\mathbb{Q}(\alpha)$ kann man jetzt auch das Polynom $X^3-2$ 
faktorisieren, indem man durch den Faktor $(X-\alpha)$ dividiert.
\index{Polynomdivision}%
Die Polynomdivision ergibt
\[
\renewcommand{\arraycolsep}{1.0pt}
\begin{array}{rcrcrcrcrcrcrcrcr}
\llap{$($}
X^3 & &           & &            &-&        2\rlap{$)$}&\phantom{(}:\phantom{)}&\llap{$($} X &-& \alpha\rlap{$)$} &\phantom{(}=\mathstrut& X^2 &+& \alpha X &+& \alpha^2 \\
X^3 &-&\alpha X^2 & &            & &         & &   & &        & &     & &          & &          \\
\cline{1-7}
    & &\alpha X^2\raisebox{4pt}{\mathstrut} & &            &-&        2& &   & &        & &     & &          & &          \\
    & &\alpha X^2 &-& \alpha^2 X & &         & &   & &        & &     & &          & &          \\
\cline{3-7}
    & &           & & \alpha^2 X\raisebox{4pt}{\mathstrut} &-&        2& &   & &        & &     & &          & &          \\
    & &           & & \alpha^2 X &-& \alpha^3& &   & &        & &     & &          & &          \\
\cline{5-7}
    & &           & &            & & \alpha^3\raisebox{4pt}{\mathstrut}\rlap{$\mathstrut-2 = 0$.}& &   & &        & &     & &          & &
\end{array}
\]
Tatsächlich ergibt Ausmultiplizieren zur Kontrolle
\begin{align}
(X-\alpha)
(X^2 + \alpha X + \alpha^2)
&=
X^3 +\alpha X^2 +\alpha^2 X
-
\alpha
X^2
-
\alpha^2 X
-
\alpha^3
\label{buch:komplex:zahlen:eqn:faktorisierung}
\\
&=
X^3-\alpha^3
\notag
\\
&=
X^3-2.
\notag
\end{align}
In der üblichen algebraischen Notation würde man $\alpha=\sqrt[3]{2}$
für die Kubikwurzel schreiben.

Wir erwarten aber, dass die Gleichung $X^3-2=0$ zwei weitere Lösungen
hat, die sich aus der Faktorisierung
\eqref{buch:komplex:zahlen:eqn:faktorisierung}
berechnen lassen sollten.
Das Polynom $X^2 +\alpha X + \alpha^2$ ist allerdings irreduzibel
als Polynom in $\mathbb{Q}(\alpha)[X]$, es definiert also eine neue
Erweiterung um eine Zahl $\beta$, die das Minimalpolynom
$b(X)=X^2+\alpha X + \alpha^2$ hat.
Die Faktorisierung ist also erst in $\mathbb{Q}(\alpha)(\beta)
=\mathbb{Q}(\alpha,\beta)$ möglich.

Die konventionelle Schreibweise für die Nullstelle des Polynoms
$b$ ergibt sich aus der Lösungsformel für quadratische Gleichungen.
Man muss
\begin{equation}
\beta
=
-\frac{\alpha}{2}
+
\!\sqrt{
\rlap{$\phantom{\bigg|}$}
\smash{\biggl(\frac{\alpha}{2}\biggr)^{\raisebox{-1pt}{$\scriptstyle2$}}-\alpha^2}
}
=
-\frac{\alpha}{2}
+
\!\sqrt{-\frac34\alpha^2}
\label{buch:komplex:zahlen:eqn:cbrtnegativ}
\end{equation}
setzen.
\end{beispiel}

Das Beispiel zeigt, wie die Lösung einer Gleichung möglicherweise
das wiederholte Hinzufügen neuer Zahlen nötig macht.
Die Notation $\sqrt[3]{2}$ für das Element $\alpha$ ist gefährlich,
da sie bereits suggeriert, dass $\alpha$ eine reelle Zahl sein
soll.
Wir haben aber nur gefordert, dass $\alpha$ eine beliebige
Nullstelle des Polynoms sein soll.
Die Darstellung \eqref{buch:komplex:zahlen:eqn:cbrtnegativ} der
anderen Nullstellen zeigt für $\alpha=\sqrt[3]{2}$, dass diese
wegen
\[
-\frac34 \!\sqrt[3]{2} < 0
\]
keine reellen Zahlen sein können.
Dieses Argument verwendet die Ordnungsrelation in den reellen
Zahlen und insbesondere die Eigenschaft, dass negative Zahlen
keine reellen Quadratwurzeln haben.
\index{Ordnungsrelation}%
Auf die Ordnungsrelation wurde in der Konstruktion von
$\mathbb{Q}(\alpha)$ nirgends Bezug genommen.
Sogar die Konstruktion von $\mathbb{Q}(\!\sqrt{2})$ hätte genauso 
gut im $-\!\sqrt{2}$ an Stelle von $\alpha$ funktionieren können.
Mit rein arithmetischen Mitteln sind die beiden Nullstellen von
$X^2-2$ nicht unterscheidbar.

\begin{definition}[algebraisches Element]
Ein Element $\alpha$ heisst \emph{algebraisch über $\mathbb{Q}$},
wenn es ein Polynom $m(X)\in\mathbb{Q}[X]$ gibt, welches $m(\alpha)=0$
erfüllt.
\index{algebraisches Element}%
\index{Element, algebraisch}%
\end{definition}

Die Definition von $\sqrt{2}$ oder $\sqrt[3]{2}$ als
algebraische Elemente über $\mathbb{Q}$ ermöglicht also,
mit diesen Elementen exakt zu rechnen, ohne dass auf
Approximationen zurückgegriffen werden muss.
Die algebraische Definition ist daher die Methode der
Wahl für ein Computeralgebrasystem.

%
% Division im Erweiterungskörper
%
\subsubsection{Division in der algebraischen Erweiterung}
Die Konstruktion von $\alpha^{-1}$ in einer algebraischen Erweiterung
$\mathbb{Q}(\alpha)$ erklärt nicht, wie man durch ein allgemeines
Element wie in \eqref{buch:komplex:zahlen:def:eqn:Qa} dividiert.
Dazu muss man die multiplikative Inverse
\[
(a_0+a_1\alpha + \dots + a_{n-1}\alpha^{n-1})^{-1}
=
b_0+b_1\alpha + \dots + b_{n-1}\alpha^{n-1}
\]
bestimmen können.
Multiplizieren wir das Produkt
\[
(a_0+a_1\alpha + \dots + a_{n-1}\alpha^{n-1})
(b_0+b_1\alpha + \dots + b_{n-1}\alpha^{n-1})
=
1
\]
aus und ersetzen wir wiederholt Potenzen $\alpha^n$ durch
$-m_0-m_1\alpha-\dots-m_{n-1}\alpha^{n-1}$, entsteht ein
Gleichungssystem für die Koeffizienten $b_k$.
\index{lineares Gleichungssystem}%
Die Formulierung und Lösung dieses Gleichungssystems ist jedoch
in Matrizenschreibweise deutlich einfacher, wir verschieben daher
die allgemeine Berechnung der multiplikativen Inversen auf
Abschnitt~\ref{buch:komplex:subsection:matrixschreibweise}.

\begin{beispiel}
Wir betrachten die Körpererweiterung $\mathbb{Q}(\alpha)$, wobei
$\alpha$ das Minimalpolynom $m(X) = X^2+X+1$ hat.
Um die multiplikative Inverse von $a+b\alpha$ zu bestimmen,
suchen wir $x+y\alpha$ derart, dass
\begin{align*}
(a+b\alpha)(x+y\alpha)&=1
\\
ax + (bx+ay)\alpha + by\alpha^2 &= 1.
\intertext{Darin ersetzen wir $\alpha^2=-\alpha-1$ und erhalten}
ax + (bx+ay)\alpha - by(1+\alpha) &= 1 \\
(ax-by) + (bx+(a-b)y)\alpha &= 1
\end{align*}
Daraus lesen wir das Gleichungssystem
\index{lineares Gleichungssystem}%
\[
\left.
\renewcommand{\arraycolsep}{2pt}
\begin{array}{rcrcr}
ax &-&     by &=& 1 \\
bx &+& (a-b)y &=& 0
\end{array}
\quad\right\}
\qquad\text{mit der Lösung}\qquad
\begin{pmatrix}x\\y\end{pmatrix}
=
\frac{1}{b^2+ab-a^2}
\begin{pmatrix}
b-a\\b
\end{pmatrix}
\]
ab.
In der üblichen Schreibweise bedeutet dies, dass
\[
(a+b\alpha)^{-1}
=
\frac{(b-a)+b\alpha}{b^2+ab-a^2}
\]
die multiplikative Inverse ist.
\end{beispiel}

%
% Die Reihenfolge der Erweiterungen
%
\subsubsection{Die Reihenfolge der Erweiterungen}
Die Quadratwurzel $\!\sqrt{2}$ ist nicht in $\mathbb{Q}$, liegt aber in
der algebraischen Erweiterung $\mathbb{Q}(\!\sqrt{2})$, die mithilfe des
über $\mathbb{Q}$ irreduziblen Polynoms $m(X)=X^2-2\in\mathbb{Q}[X]$ 
konstruiert werden kann.

$\sqrt[4]{2}$ ist ein algebraisches Element über $\mathbb{Q}$, ist
aber selbst nicht Teil dieses Körpers.
Mithilfe des über $\mathbb{Q}$ irreduziblen Polynoms $m(X)=X^4-2$.
kann man einen Körper $\mathbb{Q}(\sqrt[4]{2})$ konstruieren, welcher
das Element enthält.

Auch der Körper $\mathbb{Q}(\!\sqrt{2})$ enthält die vierte Wurzel
$\sqrt[4]{2}$ nicht.
Um einen Körper zu konstruieren, der auch $\sqrt[4]{2}$ enthält, muss
$\mathbb{Q}(\!\sqrt{2})$ um die vierte Wurzel erweitert werden.
Dazu wird das Minimalpolynom von $\sqrt[4]{2}$ benötigt, welches
irreduzibel sein muss.
Das bei der Konstruktion von $\mathbb{Q}(\sqrt[4]{2})$ verwendete
Polynom $X^4-2$ ist jedoch nicht geeignet, da es über dem aktuelle
Koeffizientenkörper $\mathbb{Q}(\!\sqrt{2})$ nicht irreduzibel ist,
sondern die Faktorzerlegung
\[
X^4 - 2 = (X + \!\sqrt{2})(X - \!\sqrt{2})
\]
hat.
Beide Faktoren haben Koeffizienten in $\mathbb{Q}(\!\sqrt{2})$,
aber nicht alle Koeffizienten sind auch in $\mathbb{Q}$.
Einen Körper, der auch $\sqrt[4]{2}$ enthält, kann man also durch
Erweiterung von $\mathbb{Q}(\!\sqrt{2})$ mithilfe des Polynoms
$X^2-\!\sqrt{2}$ erhalten.
Man könnte diesen Körper als $\mathbb{Q}(\!\sqrt{2})(\sqrt[4]{2})$
schreiben.

Umgekehrt ist die Quadratwurzel $\!\sqrt{2}$ bereits in
$\mathbb{Q}(\sqrt[4]{2})$, da $\!\sqrt{2}=(\sqrt[4]{2})^2$.
Es ist also bereits $\!\sqrt{2}\in\mathbb{Q}(\sqrt[4]{2})$.
Es folgt, dass $\mathbb{Q}(\!\sqrt{2})(\sqrt[4]{2})
=
\mathbb{Q}(\sqrt[4]{2})$.
Die Erweiterung von $\mathbb{Q}$ erst um $\!\sqrt{2}$ 
und dann um $\sqrt[4]{2}$ liefert zwei verschiedene Körper,
während die Erweiterung um $\sqrt[4]{2}$ sofort den grossen
Körper liefert.

Anders sieht es bei der Erweiterung von $\mathbb{Q}$ um $\!\sqrt{2}$
und $\!\sqrt{3}$ aus.
Als Minimalpolynom von $\!\sqrt{3}$ kann $X^2-3$ verwendet werden,
welches auch irreduzibel ist über $\mathbb{Q}(\!\sqrt{2})$.
Die Erweiterung von $\mathbb{Q}(\!\sqrt{2})$ um $\!\sqrt{3}$
liefert den Körper
\begin{align*}
\mathbb{Q}(\!\sqrt{2})(\!\sqrt{3})
&=
\Bigl\{
a+b\!\sqrt{3}
\;
\Big|
\;
a,b\in\mathbb{Q}(\!\sqrt{2})
\Bigr\}
\\
&=
\Bigl\{
a_0+a_1\!\sqrt{2}+a_2\!\sqrt{3}+a_3\!\sqrt{6}
\;
\Big|
\;
a_k\in\mathbb{Q}
\Bigr\}.
\end{align*}
Die umgekehrte Reihenfolge führt auf
\begin{align*}
\mathbb{Q}(\!\sqrt{3})(\!\sqrt{2})
&=
\Bigl\{
a+b\!\sqrt{2}
\;
\Big|
\;
a,b\in\mathbb{Q}(\!\sqrt{3})
\Bigr\}
\\
&=
\Bigl\{
b_0+b_1\!\sqrt{3}+b_2\!\sqrt{2}+b_3\!\sqrt{6}
\;
\Big|
\;
b_k\in\mathbb{Q}
\Bigr\}.
\end{align*}
Offensichtlich entsteht der gleiche Körper.
Es kommt also nicht auf die Reihenfolge an und es ist zulässig,
\[
\mathbb{Q}(\!\sqrt{2})(\!\sqrt{3})
=
\mathbb{Q}(\!\sqrt{2},\!\sqrt{3})
=
\mathbb{Q}(\!\sqrt{3},\!\sqrt{2})
=
\mathbb{Q}(\!\sqrt{3})(\!\sqrt{2})
\]
zu schreiben.

%
% Transzendente Erweiterungen
%
\subsubsection{Transzendente Erweiterungen}
Da die Menge aller rationalen Polynome abzählbar unendlich ist,
kann es nur abzählbar viele über $\mathbb{Q}$ algebraische 
Zahlen geben.
Die Menge $\mathbb{R}$ der reellen Zahlen ist aber überabzählbar
unendlich, so dass die meisten reellen Zahlen nicht algebraisch
sind.
Sie heissen transzendent.

\begin{definition}[transzendent]
Ein Element $\alpha$ heisst \emph{transzendent} über $\mathbb{Q}$,
wenn es nicht algebraisch ist.
\index{transzendent}%
\end{definition}

Die Zahlen $\pi$ und $e$ sind beide transzendent.
Der Nachweis, dass eine Zahl transzendent ist, ist oft schwierig.
Schon 1844 hat Jospeh Liouville vermutet, dass $e$ transzendent 
\index{Liouville, Joseph}%
ist, der Nachweis ist aber Charles Hermite erst 1873 gelungen
\cite{buch:hermitebeweis}.
\index{Hermite, Charles}%
Für $\pi$ ist der Beweis deutlich anspruchsvoller.
Nach einem Satz von Lindemann und Weierstrass
\index{Lindemann, Charles}%
\index{Weierstrass, Karl}%
\cite{buch:lindemann-weierstrass} sind für verschiedene
algebraische Zahlen $\alpha_1,\dots,\alpha_n$
die Zahlen $e^{\alpha_1},\dots,e^{\alpha_n}$ linear unabhängig
über $\mathbb{Q}$.
Die Zahl $0$ ist sicher algebraisch.
Wäre auch $\pi$ eine algebraische Zahl, dann wäre auch $i\pi$ algebraisch.
Nach dem Satz von Lindemann-Weierstrass müssten dann $e^0=1$ und
$e^{i\pi}$ über den algebraischen Zahlen linear unabhängig sein,
was sie wegen der eulerschen Identität
\index{eulersche Identität}%
\index{Identität, eulersch}%
\[
e^{i\pi}+1=0
\]
offensichtlich nicht sind.
Der Widerspruch zeigt, dass $\pi$ transzendent sein muss.

%
% Algebraische Erweiterungen von F_p
%
\subsubsection{Algebraische Erweiterungen von $\mathbb{F}_p$}
Die Konstruktion der reellen Zahlen mithilfe von Cauchy-Folgen ist
nur möglich, weil es auf der Menge $\mathbb{Q}$ der rationalen
Zahlen eine Ordnungsrelation gibt, die der Idee des Abstands
$|a_n-a_m|$ von Folgengliedern einen Sinn gibt.
Die Definition eines algebraischen Elementes hingegen ist auch dann 
möglich, wenn es keine solche Ordnungsrelation gibt.
Die endlichen Körper $\mathbb{F}_p$ für Primzahlen $p$ sind von
dieser Art.

\begin{definition}[endliche Körper $\mathbb{F}_p$]
\index{endliche Körper}%
\index{Korper@Körper!endlich}%
Sei $p$ eine Primzahl.
Die Menge
\[
\mathbb{F}_p
=
\{ 0, \dots , p-1 \}
\]
der ganzzahligen Reste bei Division durch $p$
heisst der \emph{endlich Körper} der Ordnung $p$.
\index{endlicher Körper}%
\index{Körper, endlich}%
\index{Fp@$\mathbb{F}_p$}%
\end{definition}

In $\mathbb{F}_p$ sind wie in $\mathbb{Q}$ alle vier Grundoperationen
wohldefiniert, insbesondere auch das reziproke Element.
Diese letzte Aussage ist gleichbedeutend damit, dass die Multiplikation
\[
x\mapsto a\cdot x
\]
mit einem Element $a\in\mathbb{F}_p\setminus\{0\}$ eine Bijektion ist.
Da dies eine lineare Abbildung ist, muss man nur prüfen, ob es ein
Element $x\ne 0$ gibt derart, dass $ax=0$ in $\mathbb{F}_p$ ist.
Dies ist gleichbedeutend damit, dass 
\[
a\cdot n \equiv 0 \mod p
\qquad\Leftrightarrow\qquad
p \mid an
\]
gilt.
Da weder $a$ noch $n$ von $p$ geteilt werden können und $p$ eine
Primzahl ist, kann es eine solche Zahl $n$ nicht geben.

Da alle ganzen Zahlen, die sich nur um ein Vielfaches von $p$
unterscheiden, das gleiche Element in $\mathbb{F}_p$ darstellen,
kann es keine mit den arithmetischen Operationen verträgelich
Ordnungsrelation geben.
Schon die Unterschiedung zwischen positiven und negativen
Zahlen ist sinnlos, da jedes Element $a$ mit $0\le a<p$ auch
die äquivalente Darstellung $a-p$ hat, welche negativ ist.
Es folgt, dass es auch nicht sinnvoll sein kann, in $\mathbb{F}_p$
von Folgen und Grenzwerten zu sprechen.
So etwas wie eine Vervollständigung von $\mathbb{F}_p$ wie
die Vervollständigung von $\mathbb{Q}$ zu $\mathbb{R}$ kann es
also nicht geben.

Trotzdem ist es sinnvoll, von algebraischen Erweiturungen zu sprechen.
Wir betrachten dazu statt Polynomen über $\mathbb{Q}$ solche über
$\mathbb{F}_p$.
Wieder setzen wir
\[
\mathbb{F}_p[X]
=
\{
X^n + a_{n-1}X^{n-1} + \dots + a_1X + a_0
\mid
a_0,\dots,a_{n-1}\in \mathbb{F}_p
\}.
\]
Wie im Falle von rationalen Koeffizienten ist ein irreduzibles
Polynom eines, welches sich nicht als Produkt zweier Polynome von
positivem Grad schreiben lässt.
Dabei kann es aber durchaus passieren, dass ein Polynom, welches
über $\mathbb{Q}$ irreduzibel war, als Polynom über $\mathbb{F}_p$
betrachtet nicht mehr irreduzibel sein muss.

\begin{beispiel}
Für $p=5$ ist
\[
(X+3)(X+2)
=
X^2 + (3+2)X + 6
=
X^2 + 1
\in
\mathbb{F}_5[X].
\]
Das Polynom $X^2+1\in\mathbb{F}_5[X]$ ist also nicht irreduzibel,
im Gegensatz zum Polynom $X^2+1\in\mathbb{Q}[X]$, welches tatsächlich
irreduzibel ist.
Dies zeigt auch, dass im Gegensatz zu $\mathbb{Q}$ oder $\mathbb{R}$,
die Gleichung $X^2+1=0$ in $\mathbb{F}_p$ eine Lösung hat.
\end{beispiel}

Wie für rationale Koeffizienten lassen sich auch für ein algebraisches
Element über $\mathbb{F}_p$ die arithmetischen Grundoperationen
definieren.

\begin{beispiel}
\label{buch:komplex:zahlen:bsp:F5}
Für $p=5$ ist $X^2-2=X^2+3\in\mathbb{F}_5[X]$ ein irreduzibles Polynom,
wie man durch Durchprobieren der möglichen Faktoren $(X-a)(X-b)$ mit
$a,b\in\mathbb{F}_5$ nachprüfen kann.
Sei jetzt $\alpha$ ein Objekt, welches $\alpha^2=2$ erfüllt.
Die Rechenoperationen sind ganz analog wie im Falle rationaler
Koeffizienten definiert.
\end{beispiel}

Die Nullstelle $\alpha$ des Polynoms $X^2-2\in\mathbb{F}_4[X]$,
die in Beispiel~\ref{buch:komplex:zahlen:bsp:F5} zu $\mathbb{F}_5$
hinzugefügt worden ist, ist eine Quadratwurzel von $2$ in $\mathbb{F}_5$,
hat aber nichts mit der vertrauten Quadratwurzeln
$\!\sqrt{5}\approx 1.414213562373$ in $\mathbb{R}$ zu tun.
Während es für reelle Nullstellen von Polynomen in $\mathbb{Q}[X]$ 
immer eine Dezimaldarstellung gibt, gibt es keine alternative Darstellung
für $\alpha$.

\begin{beispiel}
Das Polynom $m(X)=X^3-2=X^3+5\in\mathbb{F}_7[X]$ ist irreduzibel, denn
wenn es faktorisiert werden könnte, müsste notwendigerweise einer der
Faktoren den Grad 1 haben.
Ein solcher Faktor wäre also von der Form $X-a$, wobei $a\in\mathbb{F}_7$
sein müsste.
Das Polynom müsste also eine Nullstelle in $\mathbb{F}_7[X]$ haben.
Die dritten Potenzen aller Elemente von $\mathbb{F}_7$ sind allerdings
\[
\begin{tabular}{|>{$}l<{$}|>{$}c<{$}>{$}c<{$}>{$}c<{$}>{$}c<{$}>{$}c<{$}>{$}c<{$}>{$}c<{$}|}
\hline
a   & 0 & 1 & 2 & 3 & 4 & 5 & 6 \\
\hline
a^3 & 0 & 1 & 1 & 6 & 1 & 1 & 6 \\
\hline
\end{tabular}
\]
Keine dritte Potenz ergibt 2.

Fügt man $\mathbb{F}_7$ das Element $\alpha$, welches eine Nullstelle
von $m$ ist, zu $\mathbb{F}_7$ hinzu, erhält man neu die Menge
$\mathbb{F}_7(\alpha)$.
Lässt man Koeffizienten aus $\mathbb{F}_7(\alpha)$ zu, lässt sich das
Polynom $m$ wie in Beispiel~\ref{buch:komplex:zahlen:bsp:kubikwurzel}
als
\[
x^3-2
=
(x-\alpha)(\underbrace{x^2+\alpha x + \alpha^2}_{\displaystyle = b})
\]
faktorisieren.
Auch in diesem Beispiel ist eine weitere Erweiterung
$\mathbb{F}_7(\alpha,\beta)$ von $\mathbb{F}_p(\alpha)$ durch eine
Nullstelle $\beta$ des Polynoms $b\in\mathbb{F}_7(\alpha)$ nötig,
damit die anderen zwei Nullstellen von $m$ ebenfalls gefunden werden
können.
\end{beispiel}

Die Adjunktion einer Nullstelle eines irreduziblen Polynoms 
\index{Adjunktion}%
aus $\mathbb{F}_p[x]$ zu $\mathbb{F}_p$ erzeugt also wieder
die Menge 
\[
\mathbb{F}_p(\alpha)
=
\{
a_0+a_1\alpha + \dots a_{n-1}\alpha^{n-1}
\mid
a_0,\dots,a_{n-1}\in\mathbb{F}_p
\},
\]
die ein Vektorraum über $\mathbb{F}_p$ ist, mit den $n$ Monomen
\[
1, \alpha, \dots, \alpha^{n-1}
\]
als Basis.
$\mathbb{F}_p(\alpha)$ enthält somit $p^n$ Elemente.

\begin{definition}[endliche Körper]
Die Menge $\mathbb{F}_p(\alpha)$ mit $p^n$ Elementen heisst der
\emph{endliche Körper mit $p^n$} Elementen.
\index{endliche Körper}%
\index{Korper@Körper!endlich}%
\end{definition}

Fügt man einem Körper mit $p^n$ Elementen ein weiteres Element mit
einem Minimalpolynom vom Grad $l$ hinzu, entsteht eine Menge mit
mit $(p^n)^l=p^{nl}$ Elementen.
Man kann zeigen, dass alle Körper mit $p^n$ wie man sagt isomorph sind,
ganz unabhängig davon, was für ein Minimalpolynom verwendet wird
und schreibt daher unabhängig vom Minimalpolynom
$\mathbb{F}_p(\alpha)=\mathbb{F}_{p^n}$.
Die endlichen Körper spielen eine bedeutende Rolle in der Kryptographie.
Zum Beispiel verwendet die Spezifikation des AES-Algorithmus die 
\index{AES-Algorithmus}%
Arithmetik in einem Körper $\mathbb{F}_{2^8}$ mit $2^8$ Elementen,
die im Code als Bytes dargestellt werden.

%
% Die ganzen gaussschen Zahlen Z[i]
%
\subsubsection{Die ganzen gaussschen Zahlen $\mathbb{Z}[i]$}
\input{chapters/010-komplex/fig/fig-gausszahlen.tex}%
Die Erweiterung einer Zahlenmenge um ein neues Element, welches als
Nullstelle eines irreduziblen Polynoms definiert wurde, lässt sich
auch auf die Menge $\mathbb{Z}$ anwenden.
Für das Polynom $m(X) = X^2+1\in\mathbb{Z}[X]$ wird die Nullstelle
mit $i$ bezeichnet.
Die Rechenregeln folgen unmittelbar aus den Rechenregeln für
Polynome mit ganzzahligen Koeffizienten durch Anwendung der
Identität $i^2+1=0$ oder $i^2=-1$.
Es entsteht die Menge der ganzen gaussschen Zahlen gemäss der
folgenden Definition.

\begin{definition}[ganze gausssche Zahlen]
\index{ganze gausssche Zahlen}%
Die Elemente
\[
\mathbb{Z}[i]
=
\{
a+bi
\mid
a,b\in\mathbb{Z}
\}
\]
heisst die Menge der ganzen gaussschen Zahlen.
Sie bildet einen kommutativen Ring mit Eins, der den Ring $\mathbb{Z}$
der ganzen Zahlen als echten Unterring enthält.
Die Komponente $a$ heisst der \emph{Realteil} von $z=a+bi$, geschrieben
$a=\operatorname{Re}z$, $b=\operatorname{Im}z$ heisst Imaginärteil.
\end{definition}

Die ganzen gaussschen Zahlen können dargestellt werden als die
Gitterpunkt in einer Ebene, deren horizontale Achse der reelle
Zahlenstrahl ist, während die vertikale Achse die imaginären
Zahlen darstellt (Abbildung~\ref{buch:komplex:zahlen:fig:gausszahlen}).

Die Ringoperationen sind
\begin{align*}
\text{Addition:}&&
(a+bi)\pm(c+di)
&=
(a\pm c) + (b\pm d)i
\\
\text{Multiplikation:}&&
(a+bi)(c+di)
&=
ac-bd
+
(ad+bc)i.
\end{align*}
Nicht jede ganze gausssche Zahl lässt sich invertieren.
Die Inverse wird durch Erweiterung mit $a-bi$ als
\[
\frac{1}{a+bi}\cdot\frac{a-bi}{a-bi}
=
\frac{a-bi}{a^2+b^2}
=
\frac{a}{a^2+b^2}
-
\frac{b}{a^2+b^2}i
\]
gefunden.
Die Inverse existert also genau dann, wenn $a^2+b^2=1$ ist.
Dies ist nur für die Zahlen $\pm 1$ und $\pm i$ der Falle.

\begin{definition}[Betrag]
Für eine ganze gausssche Zahl $z=a+bi\in\mathbb{Z}[i]$ heisst die
Zahl $|z|$ definiert durch $|z|^2 = a^2+b^2$ der \emph{Betrag} von $z$.
\index{Betrag}%
\end{definition}

Für zwei das Produkt $z=(a+bi)(c+di)$ zweier ganzer gaussscher Zahlen gilt
\begin{align*}
|z|^2
&=
(ac-bd)^2 + (ad+bc)^2
\\
&=
a^2c^2-2abcd+b^2d^2
+
a^2d^2
+
2abcd
+
b^2c^2
\\
&=
a^2c^2 + a^2d^2 + b^2c^2 + b^2d^2
\\
&=
(a^2+b^2)(c^2+d^2)
\\
&=
|a+bi|^2 |c+di|^2
\intertext{oder}
|z|
&=
|a+bi|\cdot |c+di|.
\end{align*}
Mit anderen Worten: Wenn sich eine ganze gausssche Zahl faktorisieren
lässt, dann lassen sich auch die Beträge faktorisieren.

Tatsächlich kann man im Ring $\mathbb{Z}[i]$ der ganzen gausschen Zahlen
eine Theorie der Faktorisierung aufstellen, die ganz analog ist zur
Theorie der Primfaktoren in $\mathbb{Z}$.

